\documentclass[crop,tikz,11pt]{standalone}
\usepackage{geometricfigures}
\usepackage{amsmath,amssymb}
\usepackage{pgfplots}
\pgfplotsset{compat=1.18}
\usetikzlibrary{arrows.meta,decorations.markings,patterns,intersections,calc}

\definecolor{lib}{HTML}{3B75AF}
\definecolor{rot}{HTML}{C1701E}
\definecolor{sep}{HTML}{B03030}
\definecolor{net}{HTML}{4E9A4E}
\definecolor{up}{HTML}{7A5AA8}

% ---------------------------------------------------------------------------
% The two cases of the nesting argument, drawn. This is a schematic -- round circles
% where the trained chart has the hook of two-charts -- because the statement it
% draws is topological: only which curve is inside which, and how the areas are
% therefore ordered, carries any content.
%
% The geometry is exact all the same, so the areas can be trusted against each
% other. Both panels use one pinched pair: an outer circle of radius 2 centred
% at (0,2) and an inner one of radius 0.7 centred at (0,0.7). Both pass through
% the origin, and being internally tangent there they meet nowhere else -- which
% is the one point the two separatrix branches share, the image of the unstable
% equilibrium. Their areas are in the ratio 4 : 0.49, close to the measured
% 14.79 : 1.84.
%
% The librating curves are the boundaries of the crescent shrunk by delta: the
% outer circle at radius 2-delta and the inner one at radius 0.7+delta, which
% for delta > 0 cross at two points near the origin. The four arc endpoints
% below are those crossings, computed once and written in rather than solved for
% in TeX. delta = 0.12 and 0.30.
\newcommand{\Pxy}{(0,0)}

\begin{document}
\begin{tikzpicture}[>=Stealth,
  gam/.style={sep,very thick},
  libc/.style={lib,thick},
  rotc/.style={rot,thick},
  tag/.style={font=\scriptsize,inner sep=1.5pt},
  panelcap/.style={font=\small,align=center,anchor=north,text width=6.0cm}]

% ===========================================================================
% (a) rotating outside gamma_- : the layout of the learned chart
% ===========================================================================
\begin{scope}
  % the crescent U = image of the librating region
  \fill[lib!14!bg,even odd rule] (0,2) circle (2) (0,0.7) circle (0.7);

  % librating level sets, nested outward, exhausting U
  \draw[libc] (0.99704,0.62308) arc[start angle=-54.09, end angle=234.09, radius=1.70]
                                arc[start angle=184.41, end angle=-4.41,  radius=1.00] -- cycle;
  \draw[libc] (0.68499,0.24923) arc[start angle=-68.63, end angle=248.63, radius=1.88]
                                arc[start angle=213.34, end angle=-33.34, radius=0.82] -- cycle;

  % the two separatrix branches
  \draw[gam] (0,2)   circle (2);
  \draw[gam] (0,0.7) circle (0.7);

  % rotating level sets, outside gamma_-, nested outward
  \draw[rotc] (0,2) circle (2.17);
  \draw[rotc] (0,2) circle (2.42);

  \fill[sep] \Pxy circle (1.7pt);
  \node[tag,sep,anchor=north] at (0,-0.08) {$P$};

  \node[tag,sep,anchor=west]  at (2.08,3.30) {$\gamma_-$};
  \draw[sep!60!bg,line width=0.4pt] (0.62,0.42) -- (1.30,0.06);
  \node[tag,sep,anchor=west]  at (1.28,0.02) {$\gamma_+$};
  \node[tag,lib]              at (0,2.60) {$U$};
  \node[tag,rot,anchor=south east] at (-1.75,3.55) {rotating};

  \draw[->,rot!70!bg,line width=0.5pt] (1.62,3.86) -- (2.02,4.20);
  \node[tag,rot,anchor=south west] at (1.96,4.14) {$H$};
  \draw[->,lib!70!bg,line width=0.5pt] (-0.62,3.30) -- (-0.36,3.62);
  \node[tag,lib,anchor=south east] at (-0.40,3.58) {$H$};
\end{scope}

% ===========================================================================
% (b) rotating inside gamma_- : the same picture, the two branches exchanged
% ===========================================================================
\begin{scope}[shift={(7.0,0)}]
  \fill[lib!14!bg,even odd rule] (0,2) circle (2) (0,0.7) circle (0.7);

  \draw[libc] (0.99704,0.62308) arc[start angle=-54.09, end angle=234.09, radius=1.70]
                                arc[start angle=184.41, end angle=-4.41,  radius=1.00] -- cycle;
  \draw[libc] (0.68499,0.24923) arc[start angle=-68.63, end angle=248.63, radius=1.88]
                                arc[start angle=213.34, end angle=-33.34, radius=0.82] -- cycle;

  \draw[gam] (0,2)   circle (2);
  \draw[gam] (0,0.7) circle (0.7);

  % rotating level sets, inside gamma_-, nested inward onto whatever is left
  \fill[pattern=north east lines,pattern color=fg!35!bg] (0,0.7) circle (0.18);
  \draw[fg!45!bg,line width=0.4pt] (0,0.7) circle (0.18);
  \draw[rotc] (0,0.7) circle (0.63);
  \draw[rotc] (0,0.7) circle (0.47);
  \draw[rotc] (0,0.7) circle (0.32);

  \fill[sep] \Pxy circle (1.7pt);
  \node[tag,sep,anchor=north] at (0,-0.08) {$P$};

  \node[tag,sep,anchor=west]  at (2.08,3.30) {$\gamma_+$};
  \draw[sep!60!bg,line width=0.4pt] (0.62,0.42) -- (1.30,0.06);
  \node[tag,sep,anchor=west]  at (1.28,0.02) {$\gamma_-$};
  \node[tag,lib]              at (0,2.60) {$U$};

  \draw[rot!70!bg,line width=0.4pt] (-0.50,1.06) -- (-1.24,1.86);
  \node[tag,rot,anchor=south east] at (-1.20,1.82) {rotating};
  \draw[<-,rot!70!bg,line width=0.5pt] (-0.34,0.70) -- (-0.02,0.70);
  \node[tag,rot,anchor=east] at (-0.70,0.70) {$H$};
  \draw[fg!50!bg,line width=0.4pt] (0.14,0.84) -- (0.76,1.36);
  \node[tag,fg!60!bg,anchor=south west] at (0.72,1.32) {$K$};
\end{scope}

% ===========================================================================
\node[panelcap] at (0,-0.55)
  {\textbf{(a)} rotating family outside $\gamma_-$: every rotating loop encloses
   all of $U$, and its area grows from $\mathrm{area}(\gamma_-)$. The layout of
   the learned chart.};
\node[panelcap] at (7.0,-0.55)
  {\textbf{(b)} the branches exchanged: rotating family inside $\gamma_-$,
   nesting \emph{inward}, its area falling to $\mathrm{area}(K)\geq0$.};
\end{tikzpicture}
\end{document}
